\documentclass{article}
\usepackage[paperwidth=595pt, paperheight=841pt, margin=27.5pt, top=75pt, bottom=90pt]{geometry}
\usepackage{amsmath}
\usepackage{color,pxfonts,fix-cm}
\usepackage{latexsym}
\usepackage[T1]{fontenc}
\usepackage[utf8]{inputenc}
\usepackage{pict2e}
\usepackage{wasysym}
\usepackage[english]{babel}
\usepackage{tikz}
\usetikzlibrary{calc}
\usepackage{fancyhdr}
\usepackage{multicol}
\usepackage{hyperref}
\usepackage{framed}
\usepackage{lettrine}

% --- Color Definitions ---
\definecolor{color_101230}{rgb}{0.286275,0.14902,0.423529} % Dark Purple
\definecolor{color_283006}{rgb}{1,1,1} % White
\definecolor{color_100238}{rgb}{0.282353,0.145098,0.423529} % Slightly different Purple for text
\definecolor{color_29791}{rgb}{0,0,0} % Black
\definecolor{color_282974}{rgb}{1,0.996078,0.996078} % Off-white

% --- Hyperlink Setup ---
\hypersetup{
    colorlinks=true,
    urlcolor=color_100238,
    citecolor=color_100238,
    linkcolor=color_100238
}

\fancypagestyle{firstpage}{
    \fancyhf{} % clear header/footer

    % Main header row
    \fancyhead[L]{%
        \fontsize{13.5}{15}\selectfont\color{color_29791}Short Communication\\
        \fontsize{10}{12}\selectfont\color{color_29791}Volume 08:02, 2024
    }
    \fancyhead[R]{%
        \fontsize{13.9999}{16.8}\selectfont\color{color_100238}Journal of Animal Health and Behavioural Science
    }

    % Second header row (full-width purple bar with ISSN & Open Access)
    \fancyheadoffset[L]{0pt} % ensure full width coverage
    \fancyheadoffset[R]{0pt}
    \renewcommand{\headrulewidth}{0pt} % remove default rule
    \fancyhead[C]{%
    \begin{tikzpicture}[overlay, remember picture]
        % Purple bar (shifted lower)
        \fill[color_100238] 
            ($(current page.north west) + (0,-74pt)$) 
            rectangle 
            ($(current page.north east) + (0,-95pt)$);

        % ISSN on left
        \node[anchor=west, font=\fontsize{9}{11}\selectfont, text=white] 
            at ($(current page.north west) + (10pt,-85pt)$) {ISSN: 2952-8097};

        % Open Access on right
        \node[anchor=east, font=\fontsize{9}{11}\selectfont, text=white] 
            at ($(current page.north east) + (-10pt,-85pt)$) {Open Access};
    \end{tikzpicture}%
    }
}





% --- Header and Footer for all subsequent pages ---
\fancypagestyle{subsequent}{
    \fancyhf{}
    % Header
    \fancyhead[L]{\fontsize{8.5}{10}\selectfont\color{color_29791}James J.}
    \fancyhead[R]{\fontsize{8.5}{10}\selectfont\color{color_29791}J Anim Health Behav Sci, Volume 08:02, 2024}
    % \renewcommand{\headrule}{%
    %   \vbox to 0pt{\hbox to\headwidth{\color{color_100238}\leaders\hrule height 0.5pt\hfill}\vss}%
    % }
    % Footer
    \fancyfoot[L]{%
        \begin{tikzpicture}
            \fill[color_100238] (0,0) rectangle (76.56pt, -16.3pt);
            \node[white, font=\fontsize{8}{9.6}\selectfont] at (38.28pt, -8.15pt) {Page \thepage\ of 2};
        \end{tikzpicture}%
    }
}

% --- Apply styles ---
\pagestyle{subsequent}   % Default for all pages
\thispagestyle{firstpage} % First page special style


% --- Main Document ---
\begin{document}

% Push heading below ISSN/Open Access bar
\vspace*{1cm} % adjust this value until it clears the purple bar

\begin{center}
    {\fontsize{16}{22}\selectfont\color{color_29791}
    \textbf{Understanding Vocal Communication Patterns in Bottlenose Dolphins}\par
    }
    \vspace{0cm}
    \color{color_100238}\rule{0.97\textwidth}{0.5pt}
    \vspace{0.1cm}
\end{center}


\begin{center}
\vspace*{-0.5cm}
    \fontsize{11}{13.2}\selectfont\color{color_29791}Jhump James*\par
    \vspace{0cm}
    \fontsize{9}{10.8}\selectfont\color{color_29791}Department of Psychiatry, The University of Mostaganem, Algeria, Africa\par
    \vspace{0.1cm}
    \color{color_100238}\rule{0.6\textwidth}{1pt}
\end{center}

\begin{multicols}{2}
\fontsize{9}{10.8}\selectfont\color{color_29791}

\section*{\color{color_100238}Introduction}
\lettrine[lines=2, findent=3pt, nindent=0pt]{C}{ommunication} among marine mammals, particularly dolphins, has long fascinated researchers and enthusiasts alike. Among these creatures, the bottlenose dolphin (\textit{Tursiops truncatus}) stands out not only for its intelligence and social behavior but also for its complex vocal communication system. Studying how dolphins communicate vocally provides insights into their social dynamics, cognitive abilities and evolutionary adaptations [1]. This delves into the intricacies of vocal communication in bottlenose dolphins, exploring the types of vocalizations, their functions and the methods used to study them.

Bottlenose dolphins are highly social animals, living in groups called pods that vary in size from a few individuals to over a hundred members. Within these pods, individuals establish and maintain complex social relationships, often relying heavily on vocalizations to convey information and maintain cohesion. The study of dolphin vocalizations began in earnest in the 1960s, with pioneering work by researchers such as John C. Lilly, who used innovative techniques to record and analyze dolphin sounds. The vocal repertoire of bottlenose dolphins is extensive and diverse, consisting of both whistles and clicks. Whistles are tonal sounds produced in the nasal sacs and are often associated with individual recognition and social communication. Clicks, on the other hand, are short broadband sounds used primarily for echolocation, a technique dolphins use to navigate and locate prey in their underwater environment [2]. While these two categories broadly classify dolphin vocalizations, researchers have identified a range of variations and complexities within each type.

One of the most intriguing aspects of dolphin communication is the development of signature whistles. Each dolphin develops a unique whistle early in life, often referred to as its signature whistle. This signature whistle is used by individuals to identify themselves and to communicate identity to other members of their pod. Research has shown that dolphins can recognize and respond to the signature whistles of familiar individuals, indicating a level of individualized communication that parallels human names or identifiers. In addition to signature whistles, dolphins produce a variety of other whistles that serve different social functions [3]. Contact calls, for example, are used to maintain group cohesion when dolphins are separated, such as during foraging or other activities. These calls help individuals keep track of each other's locations and activities, ensuring that the pod remains together despite the challenges of their aquatic environment. The precise acoustic structure and usage of these contact calls vary among dolphin populations, suggesting cultural differences in communication practices.

Furthermore, whistles are not limited to social contexts; they also play a role in mate attraction and reproductive behavior. During courtship and mating, dolphins produce specific whistles that may signal readiness to mate or convey other information related to reproductive status. Understanding these reproductive vocalizations is crucial for comprehending the mating strategies and dynamics within dolphin populations, as well as their implications for population genetics and conservation efforts. In contrast to whistles, clicks serve primarily as a means of echolocation. Echolocation involves the emission of high-frequency clicks that bounce off objects in the environment, allowing dolphins to form a detailed auditory image of their surroundings. This sensory ability is critical for navigation, detecting prey and avoiding obstacles in the water. Dolphins adjust the frequency, duration and intensity of their clicks based on environmental conditions, demonstrating a high degree of control and adaptability in their echolocation abilities.

The study of dolphin clicks has provided valuable insights into the sensory physiology and acoustic ecology of these marine mammals. Researchers have used hydrophones and other underwater recording devices to capture and analyze click sequences, revealing the complex patterns and strategies dolphins employ in echolocation. By studying these click sequences, scientists can infer details about dolphin behavior, such as hunting techniques, prey preferences and social interactions.

\section*{\color{color_100238}Description}
Dolphins also exhibit a remarkable ability to communicate over long distances using their vocalizations. The propagation of sound in water differs significantly from air, presenting challenges and opportunities for underwater communication. Dolphins have evolved adaptations to maximize the effectiveness of their vocal signals, including adjusting the frequency and amplitude of their calls to enhance transmission and reception. These adaptations enable dolphins to maintain contact with individuals over distances of several kilometers, even in the vast expanses of the ocean. Social context plays a crucial role in shaping dolphin vocal behavior. Observations of wild dolphin populations have revealed that vocalizations vary depending on the activities and interactions taking place within the group [4]. For example, dolphins may produce specific vocalizations during feeding, traveling, or socializing, each serving a distinct communicative purpose. The ability to adapt vocal behavior to different social contexts highlights the sophistication and flexibility of dolphin communication systems.

Moreover, the study of dolphin vocalizations extends beyond naturalistic observations to include experimental research and controlled studies. Researchers have conducted playback experiments, where recorded dolphin vocalizations are played back to individuals or groups to observe their behavioral responses. These experiments help elucidate the function and meaning of different vocalizations, providing empirical evidence for hypotheses about dolphin communication. Advancements in technology have revolutionized the study of dolphin vocalizations, allowing researchers to capture, analyze and interpret acoustic data with unprecedented precision. Digital recording equipment, underwater cameras and sophisticated computer software have enabled detailed acoustic analyses of dolphin vocalizations, revealing intricate patterns and structures that were previously inaccessible. These technological innovations continue to expand our understanding of dolphin communication, offering new avenues for research and discovery.

In addition to their scientific significance, dolphin vocalizations hold cultural and conservation value. Indigenous cultures that have coexisted with dolphins for centuries often have rich oral traditions and stories about these marine mammals, reflecting the deep cultural connections between humans and dolphins. Furthermore, the conservation of dolphin populations relies in part on understanding their communication behaviors and how these behaviors may be impacted by human activities such as noise pollution, habitat degradation and climate change [5]. Efforts to protect and conserve dolphin populations are increasingly informed by knowledge of their vocal communication patterns and social structures. Conservation biologists collaborate with acoustic researchers to monitor dolphin populations using passive acoustic monitoring techniques, which involve deploying underwater microphones (hydrophones) to listen for and record dolphin vocalizations. These acoustic monitoring efforts provide valuable data for assessing population trends, habitat use and the impacts of human activities on dolphin behavior.

\section*{\color{color_100238}Conclusion}
In conclusion, the study of vocal communication in bottlenose dolphins offers a fascinating glimpse into the complexity of non-human communication systems. From signature whistles to echolocation clicks, dolphins employ a diverse array of vocalizations to navigate their social and ecological worlds. Advances in technology and interdisciplinary research continue to enhance our understanding of dolphin communication, shedding light on the cognitive abilities, social dynamics and evolutionary adaptations of these remarkable marine mammals. By unraveling the mysteries of dolphin vocalizations, scientists deepen their appreciation of the natural world and contribute to efforts aimed at conserving and protecting dolphin populations for future generations.

\section*{\color{color_100238}Acknowledgement}
None.
\section*{\color{color_100238}Conflict of Interest}
None.

\begin{framed}
\noindent\fontsize{8.5}{10.2}\selectfont
\textbf{How to cite this article:} James, Jhump. “Understanding Vocal Communication Patterns in Bottlenose Dolphins.” \textit{J Anim Health Behav Sci} 8 (2024): 252.
\end{framed}

\section*{\color{color_100238}References}
\fontsize{8}{9.6}\selectfont
\begin{enumerate}
    \item Picciulin, Marta, Enrico Armelloni, Raffaela Falkner and Nikolina Rako-Gospić, et al. "Characterization of the underwater noise produced by recreational and small fishing boats (<14 m) in the shallow-water of the Cres-Lošinj Natura 2000 SCI." \textit{Mar Pollut. Bull} 183 (2022): 114050.
    \item Fouda, Leila, Jessica E. Wingfield, Amber D. Fandel and Aran Garrod, et al. "Dolphins simplify their vocal calls in response to increased ambient noise." \textit{Biol Lett} 14 (2018): 20180484.
    \item Boisseau, Oliver. "Quantifying the acoustic repertoire of a population: The vocalizations of free-ranging bottlenose dolphins in Fiordland, New Zealand." \textit{J Acoust Soc Am} 117 (2005): 2318-2329.
    \item La Manna, Gabriella, Nikolina Rako‐Gospić, Gianluca Sarà and Federica Gatti, et al. "Whistle variation in Mediterranean common bottlenose dolphin: The role of geographical, anthropogenic, social and behavioral factors." \textit{Ecol Evol} 10 (2020): 1971-1987.
    \item Janik, V. M. "Food–related bray calls in wild bottlenose dolphins (Tursiops truncatus)." \textit{Proc R Soc B Boil Sci} 267 (2000): 923-927.
\end{enumerate}
\end{multicols}

\vfill
\noindent
\hspace*{0pt}\color{color_100238}\rule{\textwidth}{1pt}
\par\vspace{0.3cm} % ensures space between line and text

\fontsize{8.5}{10.2}\selectfont
\color{color_100238}\textbf{*Address for Correspondence:} 
\color{color_29791}Jhump James, Department of Psychiatry, The University of Mostaganem, Algeria, Africa, 
E-mail: \href{mailto:jhumpjames1@gmail.com}{jhumpjames1@gmail.com} \\

\color{color_100238}\textbf{Copyright:} 
\color{color_29791}© 2024 James J. This is an open-access article distributed under the terms of the Creative Commons Attribution License, which permits unrestricted use, distribution and reproduction in any medium, provided the original author and source are credited. \\

\vspace{0.1cm}
\color{color_29791}\textbf{Received:} 01 April, 2024, Manuscript No. ahbs-24-140508; 
\textbf{Editor assigned:} 02 April, 2024, PreQC No. P-140508; 
\textbf{Reviewed:} 16 April, 2024, QC No. Q-140508; 
\textbf{Revised:} 22 April, 2024, Manuscript No. R-140508; 
\textbf{Published:} 29 April, 2024, 
DOI: \href{https://doi.org/10.37421/2952-8097.2024.8.252}{10.37421/2952-8097.2024.8.252}


\end{document}